\chapter{2009年福建卷（文）}

\section{单选题}

\begin{problem}
若集合 \(A = \{x \mid x > 0\}\)，\(B = \{x \mid x < 3\}\)，则 \(A \cap B\) 等于（\quad）

\choices
    {\(\{x\mid x < 0\}\)}
    {\(\{x\mid 0 < x < 3\}\)}
    {\(\{x\mid x > 3\}\)}
    {\(\R\)}
\answer{B}

\begin{solution}
由集合定义，元素 \(x\) 属于 \(A\cap B\) 当且仅当同时满足 \(x>0\) 和 \(x<3\)，所以 \(A\cap B=\{x\mid 0<x<3\}\)，故选 B.
\end{solution}
\end{problem}

\begin{problem}
下列函数中，与函数 \( y = \frac{1}{\sqrt{x}} \) 有相同定义域的是（\quad）

\choices
    {\(f(x) = \ln x\)}
    {\(f(x) = \frac{1}{x}\)}
    {\(f(x) = |x|\)}
    {\(f(x) = \e^{x}\)}
\answer{A}

\begin{solution}
函数 \(y=\frac{1}{\sqrt{x}}\) 有意义必须满足 \(x>0\)，所以其定义域为 \((0,+\infty)\). 四个选项中，只有 \(f(x)=\ln x\) 的定义域也是 \((0,+\infty)\)，故选 A.
\end{solution}
\end{problem}

\begin{problem}
一个容量为 \(100\) 的样本，其数据的分组与各组的频数如下表

\begin{center}
\begin{tabular}{|c|c|c|c|c|c|c|c|}
\hline
组别 & \((0,10]\) & \((10,20]\) & \((20,30)\) & \((30,40)\) & \((40,50]\) & \((50,60]\) & \((60,70]\) \\
\hline
频数 & \(12\) & \(13\) & \(24\) & \(15\) & \(16\) & \(13\) & \(7\) \\
\hline
\end{tabular}
\end{center}

则样本数据落在 \( (10,40) \) 上的频率为（\quad）

\choices
    {\(0.13\)}
    {\(0.39\)}
    {\(0.52\)}
    {\(0.64\)}
\answer{C}

\begin{solution}
区间 \((10,40)\) 包含分组 \((10,20]\)、\((20,30)\) 和 \((30,40)\)，对应的频数之和为 \(13+24+15=52\). 因此样本数据落在 \((10,40)\) 上的频率为 \(\frac{52}{100}=0.52\)，故选 C.
\end{solution}
\end{problem}

\begin{problem}
若双曲线 \(\frac{x^2}{a^2} -\frac{y^2}{3^2} = 1 \) （\(a > 0\)）的离心率为2，则 \(a\) 等于（\quad）

\choices
    {\(2\)}
    {\(\sqrt{3}\)}
    {\( \frac{3}{2} \)}
    {\(1\)}
\answer{B}

\begin{solution}
双曲线的实半轴长为 \(a\)，虚半轴长为 \(3\)，焦半距满足 \(c^2=a^2+3^2=a^2+9\). 由离心率 \(e=\frac{c}{a}=2\)，得 \(c=2a\)，从而 \(4a^2=a^2+9\)，即 \(a^2=3\). 因为 \(a>0\)，所以 \(a=\sqrt{3}\)，故选 B.
\end{solution}
\end{problem}

\begin{problem}
如图，某几何体的正视图与侧视图都是边长为 1 的正方形，且体积为 \(\frac{1}{2}\). 则该几何体的俯视图可以是（\quad）

\bitmapfigure[max width=0.88\linewidth]{img/2009/fujian_liberal/q5_condition.png}

\choices
    {\choicebitmap[width=0.15\paperwidth]{img/2009/fujian_liberal/q5_optA.png}}
    {\choicebitmap[width=0.15\paperwidth]{img/2009/fujian_liberal/q5_optB.png}}
    {\choicebitmap[width=0.15\paperwidth]{img/2009/fujian_liberal/q5_optC.png}}
    {\choicebitmap[width=0.15\paperwidth]{img/2009/fujian_liberal/q5_optD.png}}
\answer{C}

\begin{solution}
正视图和侧视图都是边长为 \(1\) 的正方形，说明几何体的高为 \(1\)，其底面在两个水平方向的长度均为 \(1\). 因而俯视图作为底面时，几何体的体积等于俯视图面积乘以高，即俯视图面积应为 \(\frac{1}{2}\). 四个图形的面积分别为 \(1\)、\(\frac{\pi}{4}\)、\(\frac{1}{2}\)、\(\frac{\pi}{4}\)，只有选项 C 的面积为 \(\frac{1}{2}\). 取该直角三角形为底面并取高为 \(1\)，其正视图和侧视图均可为边长 \(1\) 的正方形，所以故选 C.
\end{solution}
\end{problem}

\begin{problem}
阅读如图所示的程序框图，运行相应的程序，输出的结果是（\quad）

\bitmapinlinefigure[0.32\linewidth][max width=0.32\linewidth]{img/2009/fujian_liberal/q6_fig1.png}

\choices
    {\(1\)}
    {\(2\)}
    {\(3\)}
    {\(4\)}
\answer{D}

\begin{solution}
程序开始时 \(S=2\)，\(n=1\). 第一次循环后 \(S=\frac{1}{1-2}=-1\)，\(n=2\)；第二次循环后 \(S=\frac{1}{1-(-1)}=\frac{1}{2}\)，\(n=3\)；第三次循环后 \(S=\frac{1}{1-\frac{1}{2}}=2\)，\(n=4\). 此时满足判断条件 \(S=2\)，程序输出 \(n=4\)，故选 D.
\end{solution}
\end{problem}

\begin{problem}
已知锐角 \(\triangle ABC\) 的面积为 \(3\sqrt{3}\)，\(BC = 4\)，\(CA = 3\)，则角 \(C\) 的大小为（\quad）

\choices
    {\(75^{\circ}\)}
    {\(60^{\circ}\)}
    {\(45^{\circ}\)}
    {\(30^{\circ}\)}
\answer{B}

\begin{solution}
由三角形面积公式，\([ABC]=\frac{1}{2}BC\cdot CA\sin C\)，所以 \(3\sqrt{3}=\frac{1}{2}\times4\times3\sin C=6\sin C\)，从而 \(\sin C=\frac{\sqrt{3}}{2}\). 由于三角形 \(ABC\) 是锐角三角形，故 \(C=60^{\circ}\)，选 B.
\end{solution}
\end{problem}

\begin{problem}
定义在 \(\R\) 上的偶函数 \(f(x)\) 的部分图象如图所示，则在 \((-2,0)\) 上，下列函数中与 \(f(x)\) 的单调性不同的是（\quad）

\bitmapfigure[max width=0.34\linewidth]{img/2009/fujian_liberal/q8_fig1.png}

\choices
    {\(y = x^{2} + 1\)}
    {\(y = |x| + 1\)}
    {\( y = \begin{cases}
2x + 1, & x \geqslant 0,\\
x^3 + 1, & x < 0.
\end{cases} \)}
    {\( y = \begin{cases}
\e^x, & x \geqslant 0,\\
\e^{-x}, & x < 0.
    \end{cases} \)}
\answer{C}

\begin{solution}
由图可知，\(f(x)\) 在 \([0,2)\) 上单调递增. 又因为 \(f\) 是偶函数，在 \((-2,0)\) 上有 \(f(x)=f(|x|)=f(-x)\)，而 \(-x\) 随 \(x\) 增大而减小，所以 \(f(x)\) 在 \((-2,0)\) 上单调递减. 选项 A 中 \(y=x^2+1\) 在该区间上的导数 \(2x<0\)，选项 B 中 \(y=|x|+1=-x+1\) 单调递减，选项 C 在该区间上为 \(y=x^3+1\)，导数 \(3x^2>0\)，单调递增，选项 D 在该区间上为 \(y=\e^{-x}\)，导数 \(-\e^{-x}<0\)，单调递减. 因此单调性不同的是 C.
\end{solution}
\end{problem}

\begin{problem}
在平面直角坐标系中，若不等式组 \( \begin{cases}
x+y-1&\geqslant0,\\
x-1&\leqslant0,\\
ax-y+1&\geqslant0,
\end{cases} \)（\(a\) 为常数）所表示的平面区域的面积等于 \(2\)，则 \(a\) 的值为（\quad）

\choices
    {\(-5\)}
    {\(1\)}
    {\(2\)}
    {\(3\)}
\answer{D}

\begin{solution}
前两条不等式给出 \(y\geqslant1-x\) 和 \(x\leqslant1\)，第三条给出 \(y\leqslant ax+1\). 要使区域有有限的正面积，必须有 \(a>-1\)，此时两条斜边在 \((0,1)\) 相交，且区域的横坐标范围为 \(0\leqslant x\leqslant1\). 在横坐标为 \(x\) 处，区域的纵向长度为 \((ax+1)-(1-x)=(a+1)x\)，所以面积为 \(\int_0^1(a+1)x\,\mathrm{d}x=\frac{a+1}{2}\). 由题意 \(\frac{a+1}{2}=2\)，得 \(a=3\). 当 \(a=-1\) 时区域退化，当 \(a<-1\) 时区域向左无界，均不符合面积为 \(2\)，故选 D.
\end{solution}
\end{problem}

\begin{problem}
设 \(m\)，\(n\) 是平面 \(\alpha\) 内的两条不同直线，\(l_{1}\)，\(l_{2}\) 是平面 \(\beta\) 内的两条相交直线，则 \(\alpha \parallel \beta\) 的一个充分而不必要条件是（\quad）

\choices
    {\(m \parallel \beta\) 且 \(l_{1} \parallel \alpha\)}
    {\(m \parallel l_{1}\) 且 \(n \parallel l_{2}\)}
    {\(m \parallel \beta\) 且 \(n \parallel \beta\)}
    {\(m \parallel \beta\) 且 \(n \parallel l_{2}\)}
\answer{B}

\begin{solution}
选项 B 中，\(m\parallel l_1\)，\(n\parallel l_2\). 因为 \(l_1\) 与 \(l_2\) 相交，所以它们不平行，从而 \(m\) 与 \(n\) 也不平行；因此 \(m,n\) 的方向可以确定平面 \(\alpha\). 同理，\(l_1,l_2\) 的方向可以确定平面 \(\beta\). 两平面中分别有两条相交直线，且对应直线互相平行，所以两平面平行，即 \(\alpha\parallel\beta\). 但平面平行并不要求所取的 \(m,n\) 必须分别平行于指定的 \(l_1,l_2\)，例如在 \(\alpha\parallel\beta\) 时另取一条不平行于 \(l_1\) 的 \(m\)，条件 B 就不成立，故该条件不是必要条件，选 B.
\end{solution}
\end{problem}

\begin{problem}
函数 \(f(x)\) 的零点与 \(g(x) = 4^x + 2x - 2\) 的零点之差的绝对值不超过 \(0.25\)，则 \(f(x)\) 可以是（\quad）

\choices
    {\(f(x) = 4x - 1\)}
    {\(f(x) = (x - 1)^{2}\)}
    {\(f(x) = \e^{x} - 1\)}
    {\(f(x) = \ln \left(x - \frac{1}{2}\right)\)}
\answer{A}

\begin{solution}
令 \(g(x)=4^x+2x-2\). 因为 \(g'(x)=4^x\ln4+2>0\)，所以 \(g\) 只有一个零点，记为 \(\xi\). 又 \(g\left(\frac14\right)=\sqrt2-\frac32<0\)，且 \(g\left(\frac13\right)=\sqrt[3]{4}-\frac43>0\)，其中 \(\sqrt2<\frac32\)，\(4>\left(\frac43\right)^3\)，所以 \(\frac14<\xi<\frac13\).

选项 A 的零点为 \(\frac14\)，且 \(\left|\xi-\frac14\right|<\frac1{12}<\frac14\). 选项 B 的零点为 \(1\)，与 \(\xi\) 的距离大于 \(\frac23\)；选项 C 的零点为 \(0\)，且 \(\xi>\frac14\)；选项 D 的零点为 \(\frac32\)，与 \(\xi\) 的距离也大于 \(\frac14\). 因此只有 A 满足条件，故选 A.
\end{solution}
\end{problem}

\begin{problem}
设 \(\bs{a}\)，\(\bs{b}\)，\(\bs{c}\) 为同一平面内具有相同起点的任意三个非零向量，且满足 \(\bs{a}\) 与 \(\bs{b}\) 不共线，\(\bs{a} \perp \bs{c}\)，\(|\bs{a}| = |\bs{c}|\)，则 \(|\bs{b} \cdot \bs{c}|\) 的值一定等于（\quad）

\choices
    {以 \(\bs{a}\)，\(\bs{b}\) 为邻边的平行四边形的面积}
    {以 \(\bs{b}\)，\(\bs{c}\) 为邻边的平行四边形的面积}
    {以 \(\bs{a}\)，\(\bs{b}\) 为两边的三角形面积}
    {以 \(\bs{b}\)，\(\bs{c}\) 为两边的三角形面积}
\answer{A}

\begin{solution}
设 \(\bs{a}\) 与 \(\bs{b}\) 的夹角为 \(\theta\). 因为 \(\bs{a}\perp\bs{c}\) 且 \(|\bs{a}|=|\bs{c}|\)，在同一平面内有 \(\left|\cos\angle(\bs{b},\bs{c})\right|=|\sin\theta|\). 因而 \(|\bs{b}\cdot\bs{c}|=|\bs{b}|\,|\bs{c}|\left|\cos\angle(\bs{b},\bs{c})\right|=|\bs{b}|\,|\bs{a}|\,|\sin\theta|\)，这正是以 \(\bs{a}\)，\(\bs{b}\) 为邻边的平行四边形面积，故选 A.
\end{solution}
\end{problem}

\section{填空题}

\begin{problem}
复数 \(\i^2 (1 + \i)\) 的实部是\fillinblank{}.
\answer{\(-1\)}

\begin{solution}
因为 \(\i^2=-1\)，所以 \(\i^2(1+\i)=-(1+\i)=-1-\i\)，其实部为 \(-1\).
\end{solution}
\end{problem}

\begin{problem}
点 \(A\) 为周长等于 \(3\) 的圆周上的一个定点. 若在该圆周上随机取一点 \(B\)，则劣弧 \(\widehat{AB}\) 的长度小于 \(1\) 的概率为\fillinblank{}.
\answer{\(\frac{2}{3}\)}

\begin{solution}
从点 \(A\) 出发，沿圆周向两个方向各取长度为 \(1\) 的弧，所得两段弧上的点都满足劣弧 \(\widehat{AB}\) 的长度小于 \(1\). 因为圆周长为 \(3\)，且 \(1<\frac32\)，这两段弧不重叠，端点是否取到不影响概率，所以有利弧长为 \(2\). 因此所求概率为 \(\frac{2}{3}\).
\end{solution}
\end{problem}

\begin{problem}
若曲线 \( f(x) = a x^2 + \ln x \) 存在垂直于 \( y \) 轴的切线，则实数 \( a \) 的取值范围是\fillinblank{}.
\answer{\((-\infty,0)\)}

\begin{solution}
曲线的定义域为 \(x>0\)，切线垂直于 \(y\) 轴即切线水平，因此需要存在 \(x>0\) 使 \(f'(x)=0\). 由 \(f'(x)=2ax+\frac1x\)，得 \(2ax^2+1=0\). 当且仅当 \(a<0\) 时，该方程有正根 \(x=\frac1{\sqrt{-2a}}\)，所以 \(a\) 的取值范围为 \((-\infty,0)\).
\end{solution}
\end{problem}

\begin{problem}
五位同学围成一圈依次循环报数，规定：

\begin{circlelist}
\item 第一位同学首次报出的数为 \(1\)，第二位同学首次报出的数也为 \(1\)，之后每位同学所报出的数都是前两位同学所报出的数之和.
\item 若报出的数为 \(3\) 的倍数，则报该次的同学需拍手一次.
\end{circlelist}

当第 \(30\) 个数被报出时，五位同学拍手的总次数为\fillinblank{}.
\answer{7}

\begin{solution}
记所报数列为 \(u_1,u_2,\ldots\)，则 \(u_1=u_2=1\)，且 \(u_n=u_{n-1}+u_{n-2}\). 对 \(3\) 取模，前八项依次为 \(1,1,2,0,2,2,1,0\)，并且第九、十项又依次为 \(1,1\)，所以由递推关系可知此后每八项重复这一周期. 因此每八个数中有两个 \(3\) 的倍数，分别对应周期中的第 \(4\) 个和第 \(8\) 个数. 前 \(30\) 个数中有三个完整周期和余下六项，所以拍手次数为 \(3\times2+1=7\).
\end{solution}
\end{problem}

\section{解答题}

\begin{problem}
等比数列 \(\{a_{n}\}\) 中，已知 \(a_{1} = 2\)，\(a_{4} = 16\).

\begin{enumerate}
\item 求数列 \( \{a_{n}\} \) 的通项公式；
\item 若 \(a_{3}\)，\(a_{5}\) 分别为等差数列 \( \{b_{n}\} \) 的第 3 项和第 5 项，试求数列 \( \{b_{n}\} \) 的通项公式及前 \(n\) 项和 \( S_{n} \).
\end{enumerate}

\begin{solution}
\begin{enumerate}
\item 设等比数列 \(\{a_n\}\) 的公比为 \(q\). 由 \(a_4=a_1q^3\) 得 \(2q^3=16\)，所以 \(q=2\)，从而数列的通项公式为 \(a_n=2\cdot2^{n-1}=2^n\).
\item 由（1）得 \(a_3=8\)，\(a_5=32\). 设等差数列 \(\{b_n\}\) 的公差为 \(d\)，则 \(b_3=8\)，\(b_5=32\)，所以 \(2d=b_5-b_3=24\)，即 \(d=12\)，进而 \(b_1=b_3-2d=8-24=-16\). 因此 \(b_n=b_1+(n-1)d=-16+12(n-1)=12n-28\)，前 \(n\) 项和为 \(S_n=\frac{n}{2}\left(2b_1+(n-1)d\right)=\frac{n}{2}\left(-32+12n-12\right)=2n(3n-11)\).
\end{enumerate}
\end{solution}
\end{problem}

\begin{problem}
袋中有大小、形状相同的红、黑球各一个，现一次有放回地随机摸取 \(3\) 次，每次摸取一个球.

\begin{enumerate}
\item 试问：一共有多少种不同的结果？ 请列出所有可能的结果；
\item 若摸到红球时得 \(2\) 分，摸到黑球时得 \(1\) 分，求 \(3\) 次摸球所得总分为 \(5\) 的概率.
\end{enumerate}

\begin{solution}
\begin{enumerate}
\item 每次摸球后都放回，且每次都有红、黑两种结果，所以按三次摸球的先后顺序记录时，共有 \(2^3=8\) 种不同结果，分别为（红，红，红）、（红，红，黑）、（红，黑，红）、（黑，红，红）、（红，黑，黑）、（黑，红，黑）、（黑，黑，红）、（黑，黑，黑）.
\item 设三次摸球中摸到红球的次数为 \(r\)，则总分为 \(2r+(3-r)=r+3\). 当总分为 \(5\) 时，\(r=2\)，故有利结果为（红，红，黑）、（红，黑，红）、（黑，红，红），共 \(3\) 种. 八种结果等可能，因此所求概率为 \(\frac{3}{8}\).
\end{enumerate}
\end{solution}
\end{problem}

\begin{problem}
已知函数 \(f(x) = \sin (\omega x + \varphi)\)，其中 \(\omega > 0\)，\(|\varphi| < \frac{\pi}{2}\).

\begin{enumerate}
\item 若 \(\cos\frac{\pi}{4}\cos\varphi-\sin\frac{3\pi}{4}\sin\varphi=0\)，求 \(\varphi\) 的值；
\item 在上一小问的条件下，若函数 \(f(x)\) 的图象的相邻两条对称轴之间的距离等于 \(\frac{\pi}{3}\)，求函数 \(f(x)\) 的解析式；并求最小正实数 \(m\)，使得函数 \(f(x)\) 的图象左平移 \(m\) 个单位所对应的函数是偶函数.
\end{enumerate}

\begin{solution}
\begin{enumerate}
\item 由题设，\(\cos\frac{\pi}{4}=\sin\frac{3\pi}{4}=\frac{\sqrt{2}}{2}\)，所以原等式化为 \(\frac{\sqrt{2}}{2}\left(\cos\varphi-\sin\varphi\right)=0\)，即 \(\tan\varphi=1\). 又因为 \(|\varphi|<\frac{\pi}{2}\)，故 \(\varphi=\frac{\pi}{4}\).
\item 正弦函数图象相邻两条对称轴之间的距离为半个周期，即 \(\frac{\pi}{\omega}\). 由 \(\frac{\pi}{\omega}=\frac{\pi}{3}\) 得 \(\omega=3\)，所以 \(f(x)=\sin\left(3x+\frac{\pi}{4}\right)\). 左平移 \(m\) 个单位后，对应函数为 \(g(x)=f(x+m)=\sin\left(3x+3m+\frac{\pi}{4}\right)\). 将其展开为 \(g(x)=\sin\left(3m+\frac{\pi}{4}\right)\cos 3x+\cos\left(3m+\frac{\pi}{4}\right)\sin 3x\). 因为 \(\cos 3x\) 是偶函数而 \(\sin 3x\) 是奇函数，所以必须有 \(\cos\left(3m+\frac{\pi}{4}\right)=0\)，即 \(3m+\frac{\pi}{4}=\frac{\pi}{2}+k\pi\)，其中 \(k\in\Z\). 因而 \(m=\frac{\pi}{12}+\frac{k\pi}{3}\). 满足 \(m>0\) 的最小值为 \(m=\frac{\pi}{12}\).
\end{enumerate}
\end{solution}
\end{problem}

\begin{problem}
如图，平行四边形 \(ABCD\) 中，\(\angle DAB = 60^{\circ}\)，\(AB = 2\)，\(AD = 4\). 将 \(\triangle CBD\) 沿 \(BD\) 折起到 \(\triangle EBD\) 的位置，使平面 \(EDB \perp\) 平面 \(ABD\).

\begin{enumerate}
\item 求证：\(AB \perp DE\)；
\item 求三棱锥 \(E-ABD\) 的侧面积.
\end{enumerate}

\bitmapfigure[max width=0.42\linewidth]{img/2009/fujian_liberal/q20_fig1.png}

\begin{solution}
\begin{enumerate}
\item 由平行四边形得 \(CD\parallel AB\)，又 \(\overrightarrow{BD}=\overrightarrow{AD}-\overrightarrow{AB}\). 因此 \(\overrightarrow{AB}\cdot\overrightarrow{BD}=\overrightarrow{AB}\cdot\overrightarrow{AD}-|\overrightarrow{AB}|^2=2\times4\cos60^{\circ}-2^2=0\)，从而 \(AB\perp BD\)，并且 \(CD\perp BD\). 折叠不改变三角形的边长，所以 \(ED=CD=AB=2\)，\(EB=CB=AD=4\). 在原平面内，\(BD^2=AB^2+AD^2-2AB\cdot AD\cos60^{\circ}=12\)，于是 \(ED^2+BD^2=4+12=16=EB^2\)，故 \(ED\perp BD\). 因为平面 \(EDB\perp\) 平面 \(ABD\)，且两平面的交线为 \(BD\)，所以 \(ED\perp\) 平面 \(ABD\)，从而 \(AB\perp DE\).

\item 由 \(AB\perp BD\) 且两平面互相垂直，得 \(AB\perp\) 平面 \(EDB\)，所以 \(AB\perp EB\). 三棱锥的三个侧面面积分别为 \([EAB]=\frac{1}{2}AB\cdot EB=4\)，\([EAD]=\frac{1}{2}AD\cdot ED=4\)，\([EBD]=\frac{1}{2}BD\cdot ED=\frac{1}{2}\times2\sqrt{3}\times2=2\sqrt{3}\). 因此三棱锥 \(E-ABD\) 的侧面积为 \(4+4+2\sqrt{3}=8+2\sqrt{3}\).
\end{enumerate}
\end{solution}
\end{problem}

\begin{problem}
已知函数 \( f(x) = \frac{1}{3} x^3 + ax^2 + bx \)，且 \( f'(-1) = 0 \).

\begin{enumerate}
\item 试用含 \(a\) 的代数式表示 \(b\)；
\item 求 \( f(x) \) 的单调区间；
\item 令 \(a = -1\)，设函数 \(f(x)\) 在 \(x_{1}\)，\(x_{2}\)（\(x_{1} < x_{2}\)）处取得极值，记点 \(M(x_{1}, f(x_{1}))\)，\(N(x_{2}, f(x_{2}))\). 证明：线段 \(MN\) 与曲线 \(f(x)\) 存在异于 \(M\)，\(N\) 的公共点.
\end{enumerate}

\begin{solution}
\begin{enumerate}
\item 由 \(f'(x)=x^2+2ax+b\) 及 \(f'(-1)=0\)，得 \(1-2a+b=0\)，所以 \(b=2a-1\).
\item 代入上式，得 \(f'(x)=x^2+2ax+2a-1=(x+1)(x+2a-1)\)，其两个零点为 \(-1\) 和 \(1-2a\). 当 \(a<1\) 时，\(-1<1-2a\)，故 \(f(x)\) 在 \((-\infty,-1)\) 和 \((1-2a,+\infty)\) 上单调递增，在 \((-1,1-2a)\) 上单调递减；当 \(a=1\) 时，\(f'(x)=(x+1)^2\geqslant0\)，故 \(f(x)\) 在 \(\R\) 上单调递增；当 \(a>1\) 时，\(1-2a<-1\)，故 \(f(x)\) 在 \((-\infty,1-2a)\) 和 \((-1,+\infty)\) 上单调递增，在 \((1-2a,-1)\) 上单调递减.
\item 当 \(a=-1\) 时，由（1）得 \(b=-3\)，于是 \(f(x)=\frac{1}{3}x^3-x^2-3x\)，且 \(f'(x)=(x+1)(x-3)\). 因此 \(x_1=-1\)，\(x_2=3\)，并且 \(M\left(-1,\frac{5}{3}\right)\)，\(N(3,-9)\). 直线 \(MN\) 的斜率为 \(\frac{-9-\frac{5}{3}}{3-(-1)}=-\frac{8}{3}\)，所以其方程为 \(y=-\frac{8}{3}x-1\). 联立曲线方程与直线方程，得 \(\frac{1}{3}x^3-x^2-3x=-\frac{8}{3}x-1\)，即 \(x^3-3x^2-x+3=(x+1)(x-1)(x-3)=0\). 除去对应 \(M\)，\(N\) 的根 \(-1\)，\(3\)，还得到 \(x=1\)，且 \(1\in(-1,3)\)，所以曲线与线段 \(MN\) 的公共点 \(\left(1,-\frac{11}{3}\right)\) 异于 \(M\)，\(N\).
\end{enumerate}
\end{solution}
\end{problem}

\begin{problem}
已知直线 \(x - 2y + 2 = 0\) 经过椭圆 \(C: \frac{x^2}{a^2} + \frac{y^2}{b^2} = 1 \)（\(a > b > 0\)）的左顶点 \(A\) 和上顶点 \(D\)，椭圆 \(C\) 的右顶点为 \(B\)，点 \(S\) 是椭圆 \(C\) 上位于 \(x\) 轴上方的动点，直线 \(AS\)，\(BS\) 与直线 \(l: x = \frac{10}{3}\) 分别交于 \(M\)，\(N\) 两点.

\begin{enumerate}
\item 求椭圆 \(C\) 的方程；
\item 求线段 \(MN\) 的长度的最小值；
\item 当线段 \(MN\) 的长度最小时，在椭圆 \(C\) 上是否存在这样的点 \(T\)，使得 \(\triangle TSB\) 的面积为 \(\frac{1}{5}\)？ 若存在，确定点 \(T\) 的个数. 若不存在，说明理由.
\end{enumerate}

\FigureLayoutDeclare{img/2009/fujian_liberal/q22_fig1.png}{12.0cm}{36bp 22bp 33bp 33bp}
\bitmapfigure[max width=0.42\linewidth]{img/2009/fujian_liberal/q22_fig1.png}

\begin{solution}
\begin{enumerate}
\item 左顶点为 \(A(-a,0)\)，上顶点为 \(D(0,b)\). 因为这两个点都在直线 \(x-2y+2=0\) 上，代入 \(A\) 得 \(-a+2=0\)，所以 \(a=2\)；代入 \(D\) 得 \(-2b+2=0\)，所以 \(b=1\). 因此椭圆 \(C\) 的方程为 \(\frac{x^2}{4}+y^2=1\)，且 \(A(-2,0)\)，\(B(2,0)\).
\item 设 \(S(x,y)\)，则 \(-2<x<2\)，\(y>0\)，且 \(x^2+4y^2=4\). 直线 \(AS\) 与 \(x=\frac{10}{3}\) 的交点纵坐标为 \(\frac{16y}{3(x+2)}\)，直线 \(BS\) 与该直线的交点纵坐标为 \(\frac{4y}{3(x-2)}\). 由于 \(x-2<0\)，前者大于后者，所以 \(MN=\frac{16y}{3(x+2)}-\frac{4y}{3(x-2)}=\frac{4y(10-3x)}{3(4-x^2)}=\frac{10-3x}{3y}\).

令 \(x=2\cos\theta\)，\(y=\sin\theta\)，其中 \(0<\theta<\pi\). 则 \(MN=\frac{2}{3}\cdot\frac{5-3\cos\theta}{\sin\theta}\). 记 \(h(\theta)=\frac{5-3\cos\theta}{\sin\theta}\)，则 \(h'(\theta)=\frac{3-5\cos\theta}{\sin^2\theta}\). 当 \(\cos\theta<\frac{3}{5}\) 时，\(h'(\theta)>0\)；当 \(\cos\theta>\frac{3}{5}\) 时，\(h'(\theta)<0\). 因此当 \(\cos\theta=\frac{3}{5}\)，即 \(\sin\theta=\frac{4}{5}\) 时，\(MN\) 取得最小值，此时 \(S\left(\frac{6}{5},\frac{4}{5}\right)\)，且 \(MN_{\min}=\frac{10-3\times\frac{6}{5}}{3\times\frac{4}{5}}=\frac{8}{3}\).
\item 在上述最小情形下，设椭圆上的点 \(T(u,v)\). 由 \(B(2,0)\)，\(S\left(\frac{6}{5},\frac{4}{5}\right)\)，得 \([\triangle TSB]=\frac{1}{2}\left|\left(-\frac{4}{5}\right)v-\frac{4}{5}(u-2)\right|=\frac{2}{5}\left|2-u-v\right|\). 令其等于 \(\frac{1}{5}\)，得到 \(u+v=\frac{3}{2}\) 或 \(u+v=\frac{5}{2}\).

对椭圆上任意点，令 \(u=2\cos t\)，\(v=\sin t\)，则 \(u+v=2\cos t+\sin t\leqslant\sqrt{5}<\frac{5}{2}\)，所以直线 \(u+v=\frac{5}{2}\) 与椭圆没有公共点. 对直线 \(u+v=\frac{3}{2}\)，代入椭圆方程得 \(\frac{u^2}{4}+\left(\frac{3}{2}-u\right)^2=1\)，即 \(5u^2-12u+5=0\)，从而 \(u=\frac{6\pm\sqrt{11}}{5}\)，对应 \(v=\frac{3\mp2\sqrt{11}}{10}\). 这给出椭圆上的两个不同点，因此存在这样的点 \(T\)，共有 \(2\) 个.
\end{enumerate}
\end{solution}
\end{problem}
